\documentclass[11pt]{article}
\usepackage[a4paper,margin=25mm]{geometry}
\usepackage{amsmath,amssymb,amsthm,booktabs,hyperref}

\newtheorem{theorem}{Theorem}[section]
\newtheorem{lemma}[theorem]{Lemma}
\newtheorem{corollary}[theorem]{Corollary}
\theoremstyle{definition}
\newtheorem{definition}[theorem]{Definition}
\theoremstyle{remark}
\newtheorem{remark}[theorem]{Remark}

\title{Weak Correspondence from Solar Eclipses:\\
Barrier Lemmas and Geometric Refutation}
\author{Shanti Draconis}
\date{2026}

\begin{document}
\maketitle

\begin{abstract}
This note formalizes a weak correspondence principle for solar eclipses.
An eclipse does not prove a Millennium Prize statement. It can, however,
refute a physical model when the predicted eclipse observables disagree with
ground truth. We state the eclipse operator, a totality criterion in terms of
umbra/penumbra geometry, and a scale-refutation lemma using the last total
solar eclipse visible in Brazil on 3 November 1994. We also record the Saros
quasi-commensurability as a periodicity lemma. All statements are framed as
barrier or falsification results only.
\end{abstract}

\section{Status discipline}

This document proves a bounded falsification statement, not a prize-level
resolution. In particular:
\begin{itemize}
  \item proving that eclipse prediction is computationally tractable does not
  prove \(P=NP\);
  \item observing a smooth solar corona does not prove global regularity for
  Navier--Stokes;
  \item periodicity in eclipse recurrences does not imply any statement about
  zeros of \(\zeta(s)\).
\end{itemize}

\section{Eclipse operator and the weak correspondence principle}

\begin{definition}[Local eclipse operator]
Fix a surface point \(x\) on Earth and a time \(t\). The observable eclipse
vector is
\[
E_x(t)=\bigl(\Delta_x(t),\rho_M(t),\rho_S(t),W_x(t),V_x(t)\bigr),
\]
where \(\Delta_x\) is the apparent angular separation between the lunar and
solar centres, \(\rho_M,\rho_S\) are the apparent angular semidiameters of the
Moon and Sun, \(W_x\) is the local umbral-width observable and \(V_x\) is the
shadow-direction observable.
\end{definition}

\begin{definition}[Totality criterion]
At \(x\) and time \(t\), totality occurs only if
\[
\rho_M(t)>\rho_S(t)
\qquad\text{and}\qquad
\Delta_x(t)<\rho_M(t)-\rho_S(t).
\]
If \(\Delta_x(t)<\rho_M(t)+\rho_S(t)\) but \(\rho_M(t)\le \rho_S(t)\), the
location is in the penumbral or annular regime rather than totality.
\end{definition}

\begin{theorem}[Weak correspondence theorem]\label{thm:weak}
Let \(M\) be a geometric model of the Sun--Moon--Earth system. If the
prediction \(E_x^{(M)}(t)\) falls outside the empirical uncertainty window of
the observed vector \(E_x^{\mathrm{obs}}(t)\) for any measured total eclipse,
then \(M\) is refuted as a model of that eclipse.
\end{theorem}

\begin{proof}
The theorem is an instance of direct model testing: a model that assigns the
wrong angular inequalities, the wrong shadow width, or the wrong shadow
direction fails on an observable required for totality. Hence it cannot be a
correct model of the event under comparison.
\end{proof}

\section{Ground truth: Brazil, 3 November 1994}

The last total solar eclipse visible in Brazil occurred on 3 November 1994 and
crossed southern Brazil before reaching greatest eclipse over the South
Atlantic. We use the published event parameters in Table~\ref{tab:truth} as
the ground-truth vector.

\begin{table}[h]
\centering
\begin{tabular}{ll}
\toprule
Parameter & Ground truth \\
\midrule
Date & 1994-11-03 \\
Type & total solar eclipse \\
Saros series & 133 \\
Greatest eclipse time & 13:39:05 UT \\
Greatest eclipse coordinates & \(35.4^\circ\mathrm{S},\,34.2^\circ\mathrm{W}\) \\
Magnitude & \(1.0535\) \\
Gamma & \(-0.3522\) \\
Maximum duration of totality & \(4\ \mathrm{min}\ 23\ \mathrm{s}\) \\
Path width at greatest eclipse & \(189\ \mathrm{km}\) \\
\bottomrule
\end{tabular}
\caption{Published parameters for the 3 November 1994 total solar eclipse.
These values are reported by the NASA/Fred Espenak eclipse bulletins and
summary tables.}
\label{tab:truth}
\end{table}

\section{Lemma 1: scale refutation from umbra geometry}

\begin{lemma}[Eclipse scale refutation]\label{lem:sphericity}
Assume the physical Moon radius remains \(R_M\approx 1737\,\mathrm{km}\) and
the solar angular semidiameter remains approximately
\(\rho_S\approx 0.266^\circ\). Consider a flat-altitude model \(H_0\) in which
the Moon is only \(h=5000\,\mathrm{km}\) above the observation plane. Then the
central totality width predicted by \(H_0\) is at least
\[
W_{H_0}\approx 2h\tan\!\bigl(\rho_M^{(H_0)}-\rho_S\bigr)
> 3000\,\mathrm{km},
\]
where
\[
\rho_M^{(H_0)}=\arctan\!\left(\frac{R_M}{h}\right)\approx 19.16^\circ.
\]
In particular, \(W_{H_0}\) is incompatible with the observed width
\(W_{\mathrm{obs}}=189\,\mathrm{km}\) from Table~\ref{tab:truth}. Therefore
\(H_0\) is refuted by direct geometric contradiction.
\end{lemma}

\begin{proof}
In the flat-altitude model, the Moon appears far larger than in the observed
sky:
\[
\rho_M^{(H_0)}=\arctan(1737/5000)\approx 19.16^\circ.
\]
Subtracting the solar semidiameter gives
\[
\rho_M^{(H_0)}-\rho_S \approx 18.89^\circ.
\]
For a point displaced transversally by distance \(r\) from the shadow centre,
the apparent centre offset is \(\arctan(r/h)\). Totality would still occur
whenever \(\arctan(r/h)<\rho_M^{(H_0)}-\rho_S\), hence
\[
r<h\tan\!\bigl(\rho_M^{(H_0)}-\rho_S\bigr).
\]
Therefore the corresponding central width satisfies
\[
W_{H_0}\approx 2h\tan(18.89^\circ)\approx 3422\,\mathrm{km}.
\]
This differs from the observed \(189\,\mathrm{km}\) by more than an order of
magnitude, so Theorem~\ref{thm:weak} refutes \(H_0\).
\end{proof}

\begin{corollary}[Weak Earth--Moon scale filter]
Any cosmological or geometric model that denies the observed Earth--Moon scale
while retaining the observed solar and lunar sizes must reproduce the eclipse
vector \(E_x(t)\) in Table~\ref{tab:truth}. If it predicts a totality width of
order \(10^3\,\mathrm{km}\) rather than \(10^2\,\mathrm{km}\), it is ruled
out by the same contradiction.
\end{corollary}

\section{Lemma 2: Saros quasi-commensurability}

\begin{lemma}[Saros periodicity bound]\label{lem:saros}
Let \(T_{\mathrm{syn}}\) denote the synodic month and
\(T_{\mathrm{dra}}\) the draconic month. Using the standard astronomical
approximations
\[
T_{\mathrm{syn}}\approx 29.53059\,\mathrm{d},
\qquad
T_{\mathrm{dra}}\approx 27.21222\,\mathrm{d},
\]
one obtains
\[
223\,T_{\mathrm{syn}}\approx 6585.3216\,\mathrm{d},
\qquad
242\,T_{\mathrm{dra}}\approx 6585.3572\,\mathrm{d},
\]
so that
\[
\Delta T_{\mathrm{Saros}}
=\bigl|223\,T_{\mathrm{syn}}-242\,T_{\mathrm{dra}}\bigr|
\approx 0.0357\,\mathrm{d}<1\,\mathrm{d}.
\]
\end{lemma}

\begin{proof}
Substitute the quoted month lengths and compute the absolute difference.
\end{proof}

\begin{remark}
Lemma~\ref{lem:saros} proves only a near-commensurability statement for eclipse
recurrence. It does not identify any spectral operator, and it does not imply
any theorem about \(\zeta(s)\).
\end{remark}

\section{Boundary with the Millennium problems}

The correct logical form is:
\begin{quote}
\emph{A well-measured eclipse can refute an incorrect physical model by
contradiction on observables. It cannot, by itself, prove an open theorem in
analysis, number theory, topology, or complexity theory.}
\end{quote}

This note therefore supports only the weak correspondence theorem and the two
barrier lemmas above. It does not claim a solution to \(P\) versus \(NP\), the
Riemann hypothesis, Navier--Stokes regularity, Yang--Mills mass gap, Hodge,
BSD, or any other Millennium Prize statement.

\section*{References}

\begin{thebibliography}{9}

\bibitem{Espenak1994}
F.~Espenak and J.~Anderson,
\emph{Total Solar Eclipse of 3 November 1994},
NASA Reference Publication 1318, NASA Goddard Space Flight Center.

\bibitem{EspenakCanon}
F.~Espenak,
\emph{Five Millennium Canon of Solar Eclipses},
NASA eclipse tables and path summaries for the 1994-11-03 event.

\end{thebibliography}

\end{document}
